
Random search arises whenever an agent must locate sparse targets with little or no sensory guidance---from animal foraging to exploration, rescue, and resource prospecting---and a central open question is which movement statistics are optimal under controlled assumptions. The L\'evy foraging hypothesis proposes that scale-free motion with power-law relocation lengths can be near-optimal in two-dimensional search, yet published conclusions remain highly model-dependent because implementation details (e.g., destructive vs.\ nondestructive targets, whether relocations are interrupted by detections, and how revisits are handled) vary across studies.

In this work we develop a transparent, reproducible 2D simulation framework for step-by-step random search on an infinite plane with Poisson-distributed targets, making modeling choices explicit while avoiding ad hoc reset rules. Within a unified framework, we compare four post-detection strategies---no-interrupt, teleport, and Viswanathan restart in both destructive and non-destructive modes---sweeping the L\'evy exponent $\mu \in [1, 3]$ (equivalently, tail exponent $\alpha = \mu - 1 \in [0, 2]$). Performance is measured primarily by the pooled unique-capture rate; mean first-passage time is analyzed for selected cases. Extensive Monte Carlo simulations across sparsity levels $\lambda/r_v \in [10, 2000]$, combined with sensitivity analyses of truncation parameters and convergence studies, yield an empirical regime map of the optimal exponent $\mu^*$. The results are consistent with $\mu^* \approx 2$ ($\alpha \approx 1$) in the sparsest tested regimes for all four strategies, show a shift toward ballistic motion ($\mu^* \to 1$) at high target density for three of four strategies, and document a trapping pathology in the non-destructive zero-jump-off implementation.
